\section{Motivation}
\label{sec:intro:motivation}
\addcontentsline{tocheb}{section}{\protect\numberline{\secnumforhebrewtoc}{מוטיבציה}}

Large language models (LLMs) are increasingly deployed as \emph{autonomous
agents} that do not merely produce text but take actions in the world: they
read and write files, query databases, push code, and send messages on a
user's behalf. The Model Context Protocol (MCP), introduced by
Anthropic in late 2024~\cite{anthropic2024mcp}, has become the de facto
standard for connecting such agents to external capabilities. An MCP
\emph{server} advertises a set of \emph{tools} (for example
\texttt{read\_file}, \texttt{delete\_file}, or \texttt{write\_query}); an MCP
\emph{client}, driven by an LLM agent, discovers those tools and invokes them
over a uniform JSON-RPC interface. This uniformity is precisely what makes MCP
powerful --- and what makes it dangerous.

The overwhelming majority of published MCP-security research studies the
protocol from the \emph{client's} point of view: how a malicious or
poisoned server can trick an agent into leaking data or executing harmful
instructions~\cite{shi2025toolhijacker,jing2025mcip}. This thesis deliberately
inverts that direction. We take the standpoint of the \textbf{server as the
protected asset} and the \textbf{agent as the threat source}. A server that
exposes a corporate filesystem, a production database, or a company Slack
workspace has no native way to decide, at the moment a call arrives, whether
\texttt{delete\_file} on a private key or an unbounded
\texttt{SELECT * FROM api\_tokens} should be allowed, throttled, or denied. An
agent that has been indirectly prompt-injected~\cite{zhan2024injecagent,
debenedetti2024agentdojo} looks, to the server, exactly like a legitimate one.
The server therefore needs a way to reason about \emph{how risky each incoming
request is} before it executes.

\section{Problem Statement}
\label{sec:intro:problem}
\addcontentsline{tocheb}{section}{\protect\numberline{\secnumforhebrewtoc}{הגדרת הבעיה}}

Traditional vulnerability-scoring systems do not answer this question. CVSS,
DREAD, and the NIST risk methodologies were designed for static software and
human-driven workflows; they score a \emph{vulnerability}, not a live tool
invocation, and they are argument-blind --- they cannot distinguish
\texttt{ls /tmp} from \texttt{rm -rf /}. Recent work confirms that when these
scores are applied to AI-specific threats they collapse toward a single value
and lose discriminative power~\cite{bahar2024validity,spring2021time}. At the
same time, the emerging MCP-specific defenses are largely \emph{binary}
detectors that label a request attack-or-benign~\cite{xing2025mcpguard,
yang2025mcpsecbench}; they do not give the server a graded, interpretable
number it can threshold and act on.

This thesis asks: \emph{Can we assign every MCP tool interaction a graduated,
interpretable risk score --- computed both statically from the tool's declared
properties and dynamically from the specific request --- that a server can use
to gate, throttle, or deny agent actions before they execute, under the
threat model in which the server is the victim and the agent is the threat?}

\section{Approach and Research Objectives}
\label{sec:intro:contributions}
\addcontentsline{tocheb}{section}{\protect\numberline{\secnumforhebrewtoc}{גישה ומטרות המחקר}}

We propose to answer this question by designing, implementing, and evaluating
\textbf{McpRisk}, a defense-oriented risk-scoring framework for MCP servers.
McpRisk decomposes the risk of a tool interaction into a small set of
interpretable factors --- the \emph{sensitivity} of the asset being touched,
the \emph{blast radius} of the operation, and the \emph{impact}
(reversibility) of the tool --- and combines them into a four-band score
(\textit{low}, \textit{medium}, \textit{high}, \textit{critical}). A
\emph{static} pipeline scores each (tool, asset) pair at design time as an
upper bound; a \emph{parameter} layer then escalates that band at request time
according to the magnitude of the actual arguments (for example, a
\texttt{SELECT} with no \texttt{LIMIT} is treated as touching every row). Every
score is produced by a local, deterministic LLM working against an explicit
rubric, cross-checked by an independent LLM judge, and --- crucially ---
validated against a panel of human and framework oracles rather than asserted.

The objectives of this research are:

\begin{enumerate}[leftmargin=*]
  \item A \textbf{server-centric threat model} for MCP
        (\Cref{sec:framework:threat_model}) that treats the agent as the
        adversary and formalizes the client~$\to$~server attack direction,
        together with a mapping from MCP tools to a 13-operation
        \emph{atomic-operation} severity taxonomy.
  \item A \textbf{factor-based scoring model}
        (\Cref{sec:framework:model}) that replaces the rigid CVSS base score
        with an interpretable product of sensitivity, blast radius, and impact,
        calibrated to a risk pyramid, together with a parameter-magnitude layer
        for request-time escalation.
  \item A \textbf{proof-of-concept implementation} (\Cref{chapter:framework})
        comprising a static scanner, a six-stage scoring pipeline, a
        parameter-rubric engine, and a call ranker, driven by a single
        reproducible local model.
  \item A \textbf{multi-oracle evaluation} (\Cref{chapter:experiments}) against
        realistic MCP deployments, a human-and-framework panel, a corpus of
        captured calls, and external attack benchmarks --- reporting where the
        approach succeeds and where it does not.
\end{enumerate}

Objectives~1--3 and an initial version of objective~4 have been realized in
preliminary work (\Cref{chapter:framework,chapter:experiments}); the remaining
evaluation and the improvements it motivates form the research plan of
\Cref{chapter:discussion_and_conclusions}.

\section{Proposal Structure}
\label{sec:intro:structure}
\addcontentsline{tocheb}{section}{\protect\numberline{\secnumforhebrewtoc}{מבנה ההצעה}}

\Cref{chapter:related_work} reviews the work surrounding the problem: the MCP
attack surface, MCP-specific defenses and attacks, the safety of tool-using
agents, guardrails and runtime gating, and risk-scoring methodologies.
\Cref{chapter:framework} presents the proposed McpRisk threat model, the scoring
model and its formulas, the atomic-operation taxonomy, and the pipeline.
\Cref{chapter:experiments} reports preliminary results from a proof-of-concept
implementation. \Cref{chapter:discussion_and_conclusions} sets out the remaining
challenges, a detailed research plan with a timeline, and the expected
contributions.
